\documentclass[aps,prd,twocolumn,superscriptaddress,nofootinbib,10pt]{revtex4-2}

\usepackage{amsmath,amssymb}
\usepackage[colorlinks=true,citecolor=blue,linkcolor=blue,urlcolor=blue]{hyperref}

\begin{document}

\title{A well-posed gap in the two-color lattice literature:\\
$T_c/\sqrt{\sigma}$ for SU(2) with $N_f=3$ light fundamental flavours}

\author{Justin Pulford}
\email{pulfordj420@gmail.com}
\affiliation{Unaffiliated}

\date{\today}

\begin{abstract}
Two-color QCD has a substantial finite-temperature literature at $N_f=2$. To our knowledge,
following a literature sweep dated 2026-07, no published determination of the transition
temperature in string-tension units, $T_c/\sqrt{\sigma}$, is available for $N_f=3$ light
fundamental Dirac flavours. We state the computation cleanly: separate chiral and deconfinement
crossovers, a chiral extrapolation, a continuum statement at least two temporal extents, and an
explicit confrontation of the staggered-rooting versus Wilson trade-off in a pseudo-real theory.
The gap is scientifically interesting independent of any particular dark-sector model: flavour
dependence of $T_c/\sqrt{\sigma}$ is well mapped in SU(3), while SU(2)'s enlarged chiral pattern
and different Goldstone count leave the saturation of that flavour dependence open. This note
reports no lattice measurement and no new numerical result; its content is the gap statement and
the requirements of a serious determination. An optional transparency remark on a pre-registered
external stake is deferred to a short final paragraph and is not the purpose of the note.
\end{abstract}

\maketitle

\section{The missing number}

We ask for a finite-temperature lattice determination of
%
\begin{equation}
  T_c/\sqrt{\sigma}
\end{equation}
%
in SU(2) gauge theory with $N_f=3$ degenerate light Dirac fermions in the fundamental
representation, with four requirements stated up front.

\begin{enumerate}
\item \emph{Both transitions separately.} The chiral crossover (peak of the chiral susceptibility)
  and the deconfinement crossover (peak of the Polyakov-loop susceptibility). In two-color QCD the
  Polyakov peak weakens as the quarks lighten, opposite to the three-color behaviour~\cite{Kaczmarek1999},
  so the splitting is itself a result.
\item \emph{A chiral extrapolation.} Ensembles at $m_{\rm PS}/\sqrt{\sigma}\lesssim 1$ with the
  quark-mass slope measured. The SU(3) benchmark is
  $T_c/\sqrt{\sigma} = 0.40(1) + 0.039(4)\,(m_{\rm PS}/\sqrt{\sigma})$~\cite{Karsch2001}; whether a
  pseudo-real theory follows the same pattern is untested.
\item \emph{A scale set by the string tension} directly, or by $r_0/t_0/w_0$ with the conversion
  stated on the same ensembles.
\item \emph{A continuum statement} from at least two temporal extents ($N_\tau=6,8$; $10$--$12$ if
  affordable). Coarse-lattice SU(3) values drifted downward by about $10\%$ in the continuum; whether
  SU(2) does the same is unknown.
\end{enumerate}

\section{Why the gap is interesting without any model}

The literature at $N_f=2$ is not empty. On the staggered side, Astrakhantsev \emph{et al.}~\cite{Astrakhantsev2019}
report $\sqrt{\sigma_0}=476(5)\,$MeV on-ensemble, and Kudrov, Bornyakov and Goy~\cite{Kudrov2025}
place the deconfinement crossover at $T_d(0)=230(10)\,$MeV, equivalently $r_0 T_d(0)=0.55(4)$.
Taken together these give $T_d/\sqrt{\sigma}\approx 0.48$. On the Wilson side, two determinations of
the chiral crossover disagree at the $\sim\!30\%$ level~\cite{Iida2020,Iida2024}.

To our knowledge (literature sweep, 2026-07), what is missing is a published
$T_c/\sqrt{\sigma}$ for $N_f=3$ light fundamentals. $N_f=4$ exists only in qualitative form from
older work. That absence matters for three reasons that do not require a dark-sector story.

\emph{Flavour dependence in a pseudo-real theory.} In SU(3) the flavour suppression of
$T_c/\sqrt{\sigma}$ largely saturates after the first two flavours. SU(2)'s fundamental
representation is pseudo-real: the chiral symmetry is enlarged, $\mathrm{SU}(2N_f)\to\mathrm{Sp}(2N_f)$,
the Goldstone count differs (fourteen at $N_f=3$), and whether the SU(3) saturation pattern
transfers is an open, clean question.

\emph{Conformal-window safety.} $N_f=3$ sits well below the SU(2) fundamental window edge near
$N_f\approx 6$~\cite{Karavirta2012,Appelquist2016}, so the theory confines and breaks chiral symmetry.
The computation is conventional.

\emph{Cost.} Two colors, three light flavours, moderate volumes; existing $N_f=2$ ensembles and
tuning are a natural starting point.

\section{The discretization is the hard choice}

$N_f=3$ is the least comfortable flavour count to discretize, and anyone costing the computation will
see that first.

\emph{Staggered.} Three flavours require the odd root, $\det^{3/4}$ --- rooted staggered where it is
least comfortable. Two-color QCD adds a second problem: with fundamental quarks the staggered action
does not reproduce the continuum symmetry-breaking pattern~\cite{Astrakhantsev2019}. Since the
pseudo-real enlargement is precisely what makes the $N_f$ dependence interesting, a discretization
that deforms it is not a neutral choice.

\emph{Wilson.} The symmetry pattern is correct, at the cost of explicit chiral-symmetry breaking at
finite lattice spacing --- which is exactly what a chiral-crossover determination is trying to
measure --- and the two existing $N_f=2$ Wilson determinations already disagree at $\sim\!30\%$.

Neither choice is settled here. The discretization is part of the computation, not a detail beneath
it. A result whose discretization is not argued for will not settle the flavour question at the
level the literature would need.

\section{Conclusion}

To our knowledge, no published $T_c/\sqrt{\sigma}$ is available for SU(2) with $N_f=3$ light
fundamental flavours (literature sweep, 2026-07). The calculation is conventional, well posed, and
discretization-sensitive. Filling the gap would answer a clean question about flavour dependence
in a pseudo-real theory, independent of any particular beyond-Standard-Model motivation. We have
stated the requirements and the literature anchors at $N_f=2$. \emph{This note reports no lattice
result}: it contains no measurement, no continuum extrapolation, and no claimed central value of
$T_c/\sqrt{\sigma}$. The computation itself remains to be done.

\paragraph{Optional transparency.}
An independent phenomenological program has publicly pre-registered a numerical stake on the
chiral transition at $T_c/\sqrt{\sigma}=0.34657$ ($=\tfrac12\ln 2$), with pre-committed scoring
rules. The stake is not a lattice result and is not the purpose of this note; it is recorded only
so that a future determination can be scored against a fixed target rather than a target chosen
afterwards. Discrimination of that stake from a nearby null would need sub-percent precision
($\sigma\lesssim 0.22\%$), well below ordinary published $1$--$3\%$ determinations. Coarser
falsification windows stated by that program are its problem, not the field's. The field's problem
is the missing $N_f=3$ number.

\begin{thebibliography}{99}

\bibitem{Kaczmarek1999}
O.~Kaczmarek, F.~Karsch, and E.~Laermann,
Thermodynamics of two-colour QCD,
Nucl.\ Phys.\ B Proc.\ Suppl.\ \textbf{73}, 441 (1999);
arXiv:hep-lat/9809059.

\bibitem{Karsch2001}
F.~Karsch, E.~Laermann, and A.~Peikert,
Quark mass and flavour dependence of the QCD phase transition,
Nucl.\ Phys.\ B \textbf{605}, 579 (2001);
arXiv:hep-lat/0012023.

\bibitem{Astrakhantsev2019}
N.~Astrakhantsev \emph{et~al.},
Lattice study of static quark-antiquark interactions in dense quark matter,
JHEP \textbf{05}, 171 (2019);
arXiv:1808.06466.

\bibitem{Kudrov2025}
I.~Kudrov, V.~Bornyakov, and V.~Goy,
Studying properties of the SU(2) QCD by lattice field theory methods,
arXiv:2511.19789.

\bibitem{Iida2020}
K.~Iida, E.~Itou, and T.-G.~Lee,
Two-colour QCD phases and the topology at low temperature and high density,
arXiv:2008.06322.

\bibitem{Iida2024}
K.~Iida, E.~Itou, and T.-G.~Lee,
JHEP \textbf{10}, 022 (2024).

\bibitem{Karavirta2012}
T.~Karavirta \emph{et~al.},
arXiv:1111.4104.

\bibitem{Appelquist2016}
T.~Appelquist \emph{et~al.},
arXiv:1511.01968.

\end{thebibliography}

\end{document}
